Sources : études citées, présentées dans le complément « Les élasticités estimées sur données françaises » ; chaque chiffre, avec sa page et son tableau, figure dans donnees/elasticites\_france.csv.\par Lecture : les deux panneaux mesurent des grandeurs différentes et ne se comparent pas. Chez Garbinti et al., les deux baisses sont estimées par la même méthode, aux seuils de tranche de 2006 à 2010, au seuil de la déclaration simplifiée entre 2012 et 2013. Chez Fack et Landais, l'élasticité est négative, puisqu'un don moins cher accroît les dons déclarés : la figure en donne la valeur absolue, et la barre claire couvre les cinq spécifications de leur tableau 2.
